% ============================================================================ %
% Encoding: UTF-8 (žluťoučký kůň úpěl ďábelšké ódy)
% ============================================================================ %

\sectionwithoutnumbering{Úvod}

První řádek prvního odstavce v kapitole či podkapitole se neodsazuje, ostatní ano. Vertikální odsazení mezi odstavci je typické pro anglickou sazbu; czech babel toto respektuje, netřeba do textu přidávat jakékoliv explicitní formátování, následuje ukázka sazby tohoto a~následujcících odstavců). Je vhodné znovu připomenout, že autor práce by si měl pohlídat, čím končí řádky. Určitě by to neměly být jednohláskové spojky ani předložky. Pro takové případy zde máme tildu (\textasciitilde), kterou \LaTeX vyhodnocuje jako nezlomitelnou mezeru. Příklad použití, ze kterého se můžete poučit, je uveden např. právě na konci tohoto řádku. Konec řádku zde a~další řádek hned za ním, aneb podívejte se do zdrojového kódu za spojku „a“.

Šablona je nastavena na oboustranný tisk. Za tímto účelem obsahuje na svém počátku místy prázdnou stranu, aby např. zadání začínalo na liché stránce (vpravo).


% ============================================================================ %
